\documentclass{../note}
\usepackage{enumitem}

\begin{document}
\begin{zadanie}{1}
\begin{enumerate}[label=(\alph*)] % <-- This syntax now works
        \item Oblicz granicę $\lim_{n\rightarrow\infty} \left(1+\frac{1}{n^{2}}\right) \left(1+\frac{2}{n^{2}}\right) \cdots \left(1+\frac{n}{n^{2}}\right).$
        
        \item Oblicz granicę $\lim_{n\rightarrow\infty} \frac{\ln n}{n} \sum_{k=2}^{n} \frac{1}{\ln k}.$
        
        \item Oblicz granicę $\lim_{n\rightarrow\infty} \left(e^{\sqrt{n+1}-\sqrt{n-1}}-1\right)^{2} \left(1+\sqrt{2}+\sqrt[3]{3}+\dots+\sqrt[n]{n}\right).$
        
        \item Oblicz granicę $\lim_{n\rightarrow\infty} \left(\sqrt{n^{2}+n}-\sqrt{n^{2}-n}\right)^{n^{2}}.$
        
        \item Oblicz granicę $\lim_{n\rightarrow\infty} \sqrt[n]{\binom{4n}{n}}.$
        
        \item Oblicz granicę $\lim_{n\rightarrow\infty} \frac{\sqrt[n^{2}]{1^{1}\cdot 2^{2}\cdot\ldots\cdot n^{n}}}{\sqrt{n}}.$
    \end{enumerate}
\end{zadanie}
\begin{rozwiazanie}
\begin{enumerate}[label=(\alph*)]
    \item Udowodnijmy następujący lemat: dla $0 \leq k \leq n - 1$ zachodzi
    \[\left(1 + \frac{1}{n^2}\right)^k \leq 1 + \frac{k+1}{n^2} \leq \left(1 + \frac{1}{n^2}\right)^{k+1}\]
    Prawa nierówność wynika wprost z nierówności Bernoulliego. Żeby udowodnić lewą, zacznijmy od następnej nierówności:
    \[k^2 + k = k(k + 1) \leq (n - 1)n \leq n^2 - n\]
    Możemy więc dodać stronami $0 \leq n^2 - n - k^2 - k$ do $n^4 \leq n^4 + n$ i otrzymać
    \[n^4 \leq n^4 - k^2 + n^2 - k\]
    Po podzieleniu stronami przez $n^2(n^2 - k)$:
    \[\frac{n^2}{n^2 - k} \leq \frac{n^2 + k + 1}{n^2}\]
    Lewa strona jest równa $\frac{1}{1 - \frac{k}{n^2}}$, czyli z odwrotnej nierówności Bernoulliego jest większa lub równa $\left(1 + \frac{1}{n^2}\right)^k$, a prawa strona jest równa $1 + \frac{k + 1}{n^2}$, co kończy dowód lematu.\\
    Możemy ograniczyć ciąg z treści od góry i od dołu:
    \[\left(1 + \frac{1}{n^2}\right)^{n^2-n} \leq \left(1+\frac{1}{n^{2}}\right) \left(1+\frac{2}{n^{2}}\right) \cdots \left(1+\frac{n}{n^{2}}\right) \leq \left(1 + \frac{1}{n^2}\right)^{n^2+n}\]
    Zarówno odgórne, jak i oddolne szacowanie różni się od $\left(1 + \frac{1}{n^2}\right)^{n^2}$ o czynnik $c = \left(1 + \frac{1}{n^2}\right)^{n}$. Możemy skorzystać z faktu, że jeśli $\lim_{n\to\infty} a_n = 0$ i $\lim_{n\to\infty} a_nb_n = g$, to $\lim_{n\to\infty} \left(1 + a_n\right)^{b_n} = e^g$. Z tego faktu wynika, że $c = 1$, czyli z tw. o trzech ciągach,
    \[\lim_{n\to\infty} \left(1+\frac{1}{n^{2}}\right) \left(1+\frac{2}{n^{2}}\right) \cdots \left(1+\frac{n}{n^{2}}\right) = \lim_{n\to\infty} \left(1 + \frac{1}{n^2}\right)^{n^2}\]
    Zatem korzystając z tego samego faktu dostajemy, że ciąg z treści zbiega do $e$.
    \item Spełnione są założenia tw. Stolza-Cezaro dla $a_n = \sum_{k=2}^{n} \frac{1}{\ln k}$, $b_n = \frac{n}{\ln n}$, czyli równie dobrze możemy obliczyć
    \[\lim_{n\to\infty} \frac{\frac{1}{\ln(n + 1)}}{\frac{n + 1}{\ln(n + 1)} - \frac{n}{\ln n}} = \lim_{n\to\infty} \frac{\ln n}{(n + 1)\ln n - n\ln(n + 1)} = \lim_{n\to\infty} \frac{1}{1 - \frac{1}{\ln n} \frac{\ln(1+\frac1n)}{\frac1n}}\]
    Wiadomo, że gdy $x \to 0$, to $\frac{\ln(1 + x)}{x} \to 1$, czyli $\frac{1}{\ln n} \frac{\ln(1+\frac1n)}{\frac1n} \to 0$, czyli powyższa granica jest równa 1.
    \item W tym podpunkcie oraz w następnym przyda się taki fakt: dla $x$ dostatecznie bliskich $0$ zachodzi 
    \[1 - \frac{x}{2} - \frac{x^2}{8} - \frac{x^3}{8} \leq \sqrt{1 - x} \leq 1 - \frac{x}{2} - \frac{x^2}{8} - \frac{x^3}{16}\]
    Można to sprawdzić podnosząc stronami do kwadratu. Zachodzi:
    \[\lim_{n\to\infty} \left(e^{\sqrt{n+1}-\sqrt{n-1}}-1\right)^{2} \left(1+\sqrt{2}+\sqrt[3]{3}+\dots+\sqrt[n]{n}\right) =\]
    \[\lim_{n\to\infty} \left(\sqrt{n+1}-\sqrt{n-1}\right)^2 \left(\frac{e^{\sqrt{n+1}-\sqrt{n-1}}-1}{\sqrt{n+1}-\sqrt{n-1}}\right)^2 \left(1+\sqrt{2}+\sqrt[3]{3}+\dots+\sqrt[n]{n}\right) =\]
    \[\lim_{n\to\infty} 2n\left(1 - \sqrt{1 - \frac{1}{n^2}}\right) \left(1+\sqrt{2}+\sqrt[3]{3}+\dots+\sqrt[n]{n}\right)\]
    Pierwszą część powyższej granicy możemy szacować zarówno odgórnie, jak i oddolnie. Zauważmy, że te szacowania są asymptotycznie równe:
    \[\lim_{n\to\infty} \frac{2n\left(\frac{1}{2n^2} + \frac{1}{8n^4} + \frac{1}{8n^6}\right)}{2n \frac{1}{2n^2}} = \lim_{n\to\infty} 1 + \frac{1}{4n^2} + \frac{1}{4n^4} = 1\]
    Z tego i z tw. o trzech ciągach wynika, że ta pierwsza część jest asymptotycznie równa swoim szacowaniom. Tak więc granica, którą liczymy, jest równa
    \[\lim_{n\to\infty} \frac1n \left(1+\sqrt{2}+\sqrt[3]{3}+\dots+\sqrt[n]{n}\right)\]
    Zauważmy, że zachodzą warunki twierdzenia Stolza-Cezaro dla \[a_n = 1+\sqrt{2}+\sqrt[3]{3}+\dots+\sqrt[n]{n}, \hspace{5pt}b_n = n\]
    Możemy więc równie dobrze obliczyć taką granicę:
    \[\lim_{n\to\infty} \frac{\sqrt[n+1]{n+1}}{n+1-n} = \lim_{n\to\infty} \sqrt[n]{n} = 1\]
    \item Zachodzi
    \[\left(\sqrt{n^{2}+n}-\sqrt{n^{2}-n}\right)^{n^{2}} = \left(2n^2\left(1 - \sqrt{1 - \frac1{n^2}}\right)\right)^{\frac{n^{2}}{2}}\]
    Możemy to oszacować zarówno odgórne, jak i oddolnie za pomocą nierówności z poprzedniego podpunktu:
    \[\left(1 + \frac{1}{4n^2} + \frac{1}{8n^4}\right)^{\frac{n^{2}}{2}} \leq \left(2n^2\left(1 - \sqrt{1 - \frac1{n^2}}\right)\right)^{\frac{n^{2}}{2}} \leq \left(1 + \frac{1}{4n^2} + \frac{1}{4n^4}\right)^{\frac{n^{2}}{2}}\]
    Teraz możemy skorzystać z faktu że gdy $a_n \to 0$ i $a_nb_n \to g$, to $(1 + a_n)^{b_n} \to e^g$. Żeby obliczyć granicę oddolnego szacowania, możemy podstawić $a_n = \frac{1}{4n^2} + \frac{1}{8n^4}$, $b_n = \frac{n^{2}}{2}$. Wtedy
    \[\lim_{n\to\infty} a_nb_n = \lim_{n\to\infty} \frac18 + \frac{1}{16n^2} = \frac18\]
    czyli granica oddolnego szacowania jest równa $e^{\frac18}$. Wychodzi że granica odgórnego szacowania też jest równa tyle, czyli z tw. o trzech ciągach, szukana granica jest równa $e^{\frac18}$.
    \item Skorzystajmy z oszacowań na silnię: $e\left(\frac{n}{e}\right)^n \leq n! \leq ne\left(\frac{n}{e}\right)^n$. Wyrażenie, którego granicę liczymy, można oszacować z obu stron:
    \[\sqrt[n]{\frac{e\left(\frac{4n}{e}\right)^{4n}}{3ne\left(\frac{3n}{e}\right)^{3n}ne\left(\frac{n}{e}\right)^n}} \leq \sqrt[n]{\frac{(4n)!}{(3n)!n!}} \leq \sqrt[n]{\frac{4ne\left(\frac{4n}{e}\right)^{4n}}{e\left(\frac{3n}{e}\right)^{3n}e\left(\frac{n}{e}\right)^n}}\]
    Zauważmy, że odgórne i oddolne szacowania różnią się o czynnik $\sqrt[n]{12n^3}$, który zbiega do 0 ponieważ $\sqrt[n]{n}$ zbiega do 0. Wychodzi, że odgórne i oddolne szacowania są asymptotycznie równe, czyli z tw. o trzech ciągach wynika, że granica, którą liczymy, jest równa granicy szacowania odgórnego.
    \[\lim_{n\to\infty} \sqrt[n]{\frac{4ne\left(\frac{4n}{e}\right)^{4n}}{e\left(\frac{3n}{e}\right)^{3n}e\left(\frac{n}{e}\right)^n}} = \lim_{n\to\infty} \frac{256}{27} \sqrt[n]{\frac{4n}{e}} = \frac{256}{27}\]
    \item Weźmy logarytm z tego wyrażenia:
    \[\log\left(\frac{\sqrt[n^{2}]{1^{1}\cdot 2^{2}\cdot\ldots\cdot n^{n}}}{\sqrt{n}}\right) = \frac{\left(\sum_{k = 1}^{n} k\log k\right) - \frac12n^2\log n}{n^2}\]
    Mianownik jest dodatni, monotoniczny i zbiega do $\infty$, czyli możemy użyć tw. Stolza-Cezaro. Musimy więc policzyć taką granicę:
    \[\lim_{n\to\infty} \frac{(n+1)\log(n+1) - \frac12(n+1)^2\log(n+1) + \frac12n^2\log n}{(n+1)^2 - n^2} =\]
    \[\lim_{n\to\infty} -\frac{1}{4+\frac2n} \frac{\log(1 + \frac1n)}{\frac1n} + \frac{\log(n+1)}{4n+2} = -\frac14 \cdot 1 + 0 = -\frac14\]
    Drugi składnik w powyższym wyrażeniu zbiega do 0 z hierarchii ciągów. Skorzystaliśmy też z faktu, że dla $x \to 0$ zachodzi $\frac{\log(1+x)}{x} \to 1$. Wychodzi, że początkowa granica to $e^{-\frac14}$.
\end{enumerate}
\end{rozwiazanie}

\begin{zadanie}{2}
Zbadaj, w zależności o parametru \( a_{0}\in\mathbb{R} \), zbieżność ciągu zadanego rekurencyjnie
\[ a_{n+1}=a_{n}^{3}-6a_{n}^{2}+12a_{n}-6, \quad n\ge0. \]
Jeśli ciąg jest zbieżny, oblicz jego granicę.
\end{zadanie}
\begin{rozwiazanie}
Zauważmy, że $a_{n+1} - 2 = (a_n - 2)^3$, czyli $a_n - 2 = (a_0 - 2)^{3^n}$. Jeśli $a_0 \in (1, 3)$, to $|a_0 - 2| < 1$, czyli $(a_0 - 2)^{3^n} \to 0$, czyli $a_n \to 2$. Jeśli $a_0 = 1$ lub $a_0 = 3$, to z indukcji wynika, że ciąg $(a_n)$ jest stały, bo $a_{n+1} = (a_n - 2)^3 + 2 = a_n$. Wtedy ciąg zbiega odpowiednio do 1 albo 3. Gdy $a_0 \notin [1, 3]$, to $|a_0 - 2| > 0$, czyli $(a_0 - 2)^{3^n}$ jest rozbieżne, czyli $a_n$ też.
\end{rozwiazanie}

\begin{zadanie}{5}
Załóżmy, że dla każdego \( \gamma>1 \) ciąg \( (a_{[\gamma^{n}]})_{n\ge1} \) jest zbieżny do 0. Czy z tego wynika, że ciąg \( (a_{n})_{n\geq1} \) jest zbieżny do 0?
\end{zadanie}
\begin{rozwiazanie}
Załóżmy nie wprost, że ciąg $(a_n)$ nie jest zbieżny do 0, czyli że istnieje nieskończony podciąg, dla którego moduły wyrazów są większe od pewnego $M > 0$. Niech ciąg $(i_n)$ to indeksy tego podciągu. Możemy skonstruować ciąg przedziałów $([l_n, r_n))$ w taki sposób, że gdy $\log(\gamma) \in [l_k, r_k)$, to $(a_{[\gamma^n]})$ zawiera $k$ elementów tego podciągu. Niech $[l_1, r_1) = [log(i_1), log(i_1 + 1))$.\\ Załóżmy, że już mamy skonstruowane $[l_k, r_k)$. Niech $c > 0$ to pewna stała oraz $n_0 = \lceil\frac{l_k + c}{r_k - l_k}\rceil$. Wtedy dla $n \geq n_0$ kolejne przedziały typu $[nl_k, nr_k)$ nachodzą na siebie o przynajmniej $c$, ponieważ ich początki są co $l_k$, a długości co najmniej równe $l_k + c$. To znaczy, że wtedy każdy przedział długości co najwyżej $c$ położony za $n_0l_k$ jest zawarty w przedziale typu $[nl_k, nr_k)$.\\
Niech $i$ to pierwszy indeks z ciągu $(b_n)$ spełniający $i \geq n_0l_k$ i nie będący jednym z $k$ indeksów, które już wybraliśmy. Wybierzmy $c$ równe długości $[log(i), log(i) + 1)$. Musi istnieć takie $n$, że przedział $[log(i), log(i) + 1)$ jest zawarty w przedziale $[nl_k, nr_k)$. Niech $[l_{k+1}, r_{k+1}) = [\frac{log(i)}{n}, \frac{log(i) + 1}{n})$. Wtedy $[l_{k+1}, r_{k+1}) \subset [l_k, r_k)$, czyli gdy $\log(\gamma) \in [l_{k+1}, r_{k+1})$, to $([\gamma^n])$ zawiera nie tylko te $k$ indeksów, które już wybraliśmy, ale też $i$, ponieważ $\gamma^n \in [i, i + 1)$.\\
Ciąg $(l_k)$ jest rosnący i ograniczony przez $r_1$, czyli zbiega do pewnej granicy $\gamma$. Dla każdego $k$ zachodzi $l_k \leq r_k$, czyli $\gamma = \sup l_k \leq r_k$, czyli $\gamma \in [l_k, r_k)$. To znaczy, że dla każdego $k$, ciąg $([\gamma^n])$ zawiera przynajmniej $k$ elementów ciągu $(b_n)$, czyli ciąg $(a_{[\gamma^n]})$ musi zawierać nieskończenie wiele elementów o module większym od $M$, co jest sprzecznością.
\end{rozwiazanie}

\end{document}